\begin{definition}[Lower ramification groups]\label{mod:ramification-lowergroups}
Let $K/F$ be a finite Galois valuative extension of nonarchimedean local fields
and write $G=\operatorname{Gal}(K/F)$ for its Galois group.  Every
$\sigma\in G$ preserves depth congruence:
\[
 \sigma(x)\equiv\sigma(y)\pmod{\mathfrak p_K^n}
 \quad\Longleftrightarrow\quad
 x\equiv y\pmod{\mathfrak p_K^n}.
\]
Consequently
\[
 C_n(a)=\{\sigma\in G:\sigma(a)\equiv a\pmod{\mathfrak p_K^n}\}
\]
is a subgroup of $G$.  Restriction to $\mathcal O_K$ and reduction give
homomorphisms
\[
 G\longrightarrow\operatorname{RingAut}(\mathcal O_K)
 \longrightarrow\operatorname{RingAut}(k_K),
\]
and the residue action sends $\overline{x}$ to $\overline{\sigma(x)}$.

For every $i\in\mathbf Z$, define
\[
 G_i=\bigcap_{x\in\mathcal O_K}C_{i+1}(x).
\]
This extends the manuscript's $i\geq-1$ convention to all integers.  The
following three membership descriptions are exactly equivalent:
\[
 \begin{aligned}
 \sigma\in G_i
 &\Longleftrightarrow
   \forall x\in\mathcal O_K,\quad
   \sigma(x)\equiv x\pmod{\mathfrak p_K^{i+1}}\\
 &\Longleftrightarrow
   \forall x\in\mathcal O_K,\quad
   i+1\leq\operatorname{ord}_K(\sigma(x)-x)\\
 &\Longleftrightarrow
   \forall x\in\mathcal O_K,\quad
   \sigma(x)-x\in\mathfrak p_K^{i+1}.
 \end{aligned}
\]
The zeroth group is the kernel of the residue action.  The filtration is
decreasing, every $G_i$ is normal in $G$, and
\[
 i\leq j\Longrightarrow G_j\leq G_i,
 \qquad i\leq-1\Longrightarrow G_i=G,
 \qquad G_{-1}=G.
\]
There is an integer $N\geq0$ such that $G_N=\{1\}$, and every $G_j$ with
$j\geq N$ is also trivial.  This is proved directly from finiteness of the
automorphism group and the fact that every nonidentity automorphism moves an
integral element; no higher ramification theory is assumed.

For every integral test element $x$, membership in $G_i$ implies the
corresponding depth-$(i+1)$ congruence on $x$.  Thus, for an arbitrary
uniformizer $\pi$,
\[
 \sigma\in G_i\Longrightarrow
 \operatorname{ord}_K(\sigma(\pi)-\pi)\geq i+1.
\]
Only this implication is asserted in general.  If an integral element $\pi$
is supplied together with the explicit generator hypothesis
\[
 \operatorname{Algebra.adjoin}_{\mathcal O_F}\{\pi\}=\mathcal O_K,
\]
then testing $\pi$ is sufficient:
\[
 \begin{aligned}
 \sigma\in G_i
 &\Longleftrightarrow
 \sigma(\pi)\equiv\pi\pmod{\mathfrak p_K^{i+1}}\\
 &\Longleftrightarrow
 \operatorname{ord}_K(\sigma(\pi)-\pi)\geq i+1,
 \end{aligned}
\]
and the exact-shell criterion is
\[
 \sigma\in G_i\setminus G_{i+1}
 \Longleftrightarrow
 \operatorname{ord}_K(\sigma(\pi)-\pi)=i+1.
\]
This is the interface used when a totally ramified extension supplies a
uniformizer satisfying the generator hypothesis.  No converse is claimed for
a general uniformizer: without such a hypothesis it is false, for example for
a nontrivial unramified automorphism fixing a base uniformizer.

For $\sigma\in G_i$ and a chosen uniformizer $\pi$, the displacement and its
normalization are retained as actual lattice elements
\[
 d_{i,\pi}(\sigma)=\sigma(\pi)-\pi\in\mathfrak p_K^{i+1},
 \qquad
 z_{i,\pi}(\sigma)=
 \frac{\sigma(\pi)-\pi}{\pi^{i+1}}\in\mathfrak p_K^0.
\]
Their coercions to $K$ are the displayed elements.  The positive-depth
residual coordinate is the quotient class
\[
 \overline d_{i,\pi}(\sigma)
 =\overline{z_{i,\pi}(\sigma)}\in k_K.
\]
Under the generator hypothesis its kernel set is exactly the next group:
\[
 \overline d_{i,\pi}(\sigma)=0
 \Longleftrightarrow \sigma\in G_{i+1}.
\]
In particular, an exact-shell element has normalized order zero and nonzero
residual displacement.

The distinct tame coordinate is also exposed.  For $\sigma\in G_0$,
\[
 r_\pi(\sigma)=\frac{\sigma(\pi)}{\pi}\in\mathfrak p_K^0,
 \qquad \operatorname{ord}_K(r_\pi(\sigma))=0,
\]
so its residue is nonzero and is bundled as
\[
 \overline r_\pi(\sigma)\in k_K^\times.
\]
Under the same generator hypothesis,
\[
 \overline r_\pi(\sigma)=1\Longleftrightarrow\sigma\in G_1
\]
both in $k_K$ and in $k_K^\times$.  These value and kernel-set declarations
are the interfaces from which the later graded-norm node bundles the tame
multiplicative and wild additive ramification maps.  The chosen uniformizer is
part of both coordinates; no coordinate-independence claim or representative
choice in $K$ is made.

For $i\leq j$, let $G_j^{(i)}$ be the range of the certified inclusion
$G_j\hookrightarrow G_i$.  Membership in $G_j^{(i)}$ is exactly membership of
the underlying automorphism in $G_j$, its image back in $G$ is $G_j$, and it
is normal in $G_i$.  Hence
\[
 \operatorname{LowerRamificationQuotient}(i,j)=G_i/G_j^{(i)}
\]
is an honest finite quotient group.  It has a specified noncomputable finite
enumeration and
\[
 \operatorname{Nat.card}(G_i/G_j^{(i)})=[G_i:G_j^{(i)}].
\]
For its canonical surjection $q_{ij}$,
\[
 \ker(q_{ij})=G_j^{(i)},\qquad
 q_{ij}(\sigma)=1\Longleftrightarrow\sigma\in G_j,
\]
and
\[
 q_{ij}(\sigma)=q_{ij}(\tau)
 \Longleftrightarrow \sigma\tau^{-1}\in G_j.
\]
The successive quotient $G_i/G_{i+1}$ and its canonical projection are
provided with the corresponding kernel-membership characterization.  No
quotient representative is selected.

Finally, for every real $u$, define lower numbering by the manuscript's
ceiling convention
\[
 G_u^{\mathbf R}=G_{\lceil u\rceil}.
\]
The real filtration is decreasing, agrees with integer lower numbering on
integer inputs, has the same exact membership criterion at depth
$\lceil u\rceil+1$, and satisfies $G_{-1}^{\mathbf R}=G$.  The definition is
available for all real $u$, although the manuscript uses it for $u\geq-1$.

The auxiliary order interfaces include
\[
 (n\leq a\ \wedge\ \neg(n+1\leq a))\Longleftrightarrow a=n
 \quad(a\in\operatorname{WithTop}(\mathbf Z))
\]
and $\operatorname{ord}_K(\pi^m)=m$ for every integer $m$ and every
uniformizer $\pi$.

\LeanFile{LanglandsFirstMainLemma/Ramification/LowerGroups.lean}
\lean{LanglandsFirstMainLemma.withTopInt_le_and_not_succ_le_iff_eq}
\lean{LanglandsFirstMainLemma.galois_congruentAtDepth_iff}
\lean{LanglandsFirstMainLemma.galoisCongruenceSubgroup}
\lean{LanglandsFirstMainLemma.mem_galoisCongruenceSubgroup}
\lean{LanglandsFirstMainLemma.galoisIntegerAut}
\lean{LanglandsFirstMainLemma.galoisResidueAction}
\lean{LanglandsFirstMainLemma.galoisResidueAction_residue}
\lean{LanglandsFirstMainLemma.lowerRamificationGroup}
\lean{LanglandsFirstMainLemma.mem_lowerRamificationGroup}
\lean{LanglandsFirstMainLemma.mem_lowerRamificationGroup_iff_ord}
\lean{LanglandsFirstMainLemma.mem_lowerRamificationGroup_iff_lattice}
\lean{LanglandsFirstMainLemma.lowerRamificationGroup_zero_eq_ker_galoisResidueAction}
\lean{LanglandsFirstMainLemma.lowerRamificationGroup_antitone}
\lean{LanglandsFirstMainLemma.lowerRamificationGroup_mono}
\lean{LanglandsFirstMainLemma.lowerRamificationGroup_succ_le}
\lean{LanglandsFirstMainLemma.lowerRamificationGroup_eq_top_of_le_neg_one}
\lean{LanglandsFirstMainLemma.lowerRamificationGroup_neg_one}
\lean{LanglandsFirstMainLemma.exists_nonnegative_lowerRamificationGroup_eq_bot}
\lean{LanglandsFirstMainLemma.lowerRamificationGroup_eq_bot_of_ge}
\lean{LanglandsFirstMainLemma.lowerRamificationGroup_eventually_bot}
\lean{LanglandsFirstMainLemma.lowerRamificationGroup_normal}
\lean{LanglandsFirstMainLemma.lowerRamificationGroup.instNormal}
\lean{LanglandsFirstMainLemma.lowerRamificationGroup_le_galoisCongruenceSubgroup}
\lean{LanglandsFirstMainLemma.lowerRamificationGroup_uniformizer}
\lean{LanglandsFirstMainLemma.lowerRamificationGroup_uniformizer_ord}
\lean{LanglandsFirstMainLemma.lowerRamificationDisplacement}
\lean{LanglandsFirstMainLemma.coe_lowerRamificationDisplacement}
\lean{LanglandsFirstMainLemma.lowerRamificationNormalizedDisplacement}
\lean{LanglandsFirstMainLemma.coe_lowerRamificationNormalizedDisplacement}
\lean{LanglandsFirstMainLemma.lowerRamificationTameRatio}
\lean{LanglandsFirstMainLemma.coe_lowerRamificationTameRatio}
\lean{LanglandsFirstMainLemma.ord_lowerRamificationTameRatio_eq_zero}
\lean{LanglandsFirstMainLemma.lowerRamificationTameResidue}
\lean{LanglandsFirstMainLemma.lowerRamificationTameResidue_ne_zero}
\lean{LanglandsFirstMainLemma.lowerRamificationTameResidueUnit}
\lean{LanglandsFirstMainLemma.coe_lowerRamificationTameResidueUnit}
\lean{LanglandsFirstMainLemma.ord_uniformizer_zpow}
\lean{LanglandsFirstMainLemma.mem_lowerRamificationGroup_iff_of_adjoin_eq_top}
\lean{LanglandsFirstMainLemma.mem_lowerRamificationGroup_iff_ord_of_adjoin_eq_top}
\lean{LanglandsFirstMainLemma.lowerRamificationGroup_mem_and_not_mem_succ_iff_ord_sub_eq}
\lean{LanglandsFirstMainLemma.ord_lowerRamificationNormalizedDisplacement_eq_zero}
\lean{LanglandsFirstMainLemma.lowerRamificationResidueDisplacement}
\lean{LanglandsFirstMainLemma.lowerRamificationResidueDisplacement_ne_zero}
\lean{LanglandsFirstMainLemma.lowerRamificationResidueDisplacement_eq_zero_iff}
\lean{LanglandsFirstMainLemma.lowerRamificationTameResidue_eq_one_iff}
\lean{LanglandsFirstMainLemma.lowerRamificationTameResidueUnit_eq_one_iff}
\lean{LanglandsFirstMainLemma.lowerRamificationGroupInside}
\lean{LanglandsFirstMainLemma.mem_lowerRamificationGroupInside}
\lean{LanglandsFirstMainLemma.map_lowerRamificationGroupInside_subtype}
\lean{LanglandsFirstMainLemma.lowerRamificationGroupInside.instNormal}
\lean{LanglandsFirstMainLemma.LowerRamificationQuotient}
\lean{LanglandsFirstMainLemma.lowerRamificationQuotientFintype}
\lean{LanglandsFirstMainLemma.lowerRamificationQuotient_natCard}
\lean{LanglandsFirstMainLemma.lowerRamificationQuotientMk}
\lean{LanglandsFirstMainLemma.lowerRamificationQuotientMk_apply}
\lean{LanglandsFirstMainLemma.lowerRamificationQuotientMk_ker}
\lean{LanglandsFirstMainLemma.lowerRamificationQuotientMk_eq_one_iff}
\lean{LanglandsFirstMainLemma.lowerRamificationQuotientMk_surjective}
\lean{LanglandsFirstMainLemma.lowerRamificationQuotientMk_eq_iff}
\lean{LanglandsFirstMainLemma.LowerRamificationGraded}
\lean{LanglandsFirstMainLemma.lowerRamificationGradedMk}
\lean{LanglandsFirstMainLemma.lowerRamificationGradedMk_eq_one_iff}
\lean{LanglandsFirstMainLemma.lowerRamificationGroupReal}
\lean{LanglandsFirstMainLemma.mem_lowerRamificationGroupReal}
\lean{LanglandsFirstMainLemma.lowerRamificationGroupReal_antitone}
\lean{LanglandsFirstMainLemma.lowerRamificationGroupReal_intCast}
\lean{LanglandsFirstMainLemma.lowerRamificationGroupReal_neg_one}
\leanok
\SourceText{Section Ramification notation and the cyclic-prime norm theorem.}
\uses{mod:localfield-extension,mod:localfield-valuation}
\FormalizationContract The construction uses only valuation congruences,
the completed extension, valuation, lattice, and residue-field APIs, and
Mathlib's elementary Galois theory.  It imports no higher-ramification theorem
or typeclass.  The converse uniformizer criterion always retains the
explicit $\mathcal O_F$-algebra generation hypothesis.  The tame and
positive-depth coordinates remain distinct; only their values and exact kernel
sets are established here, while the downstream graded-norm node bundles and
analyzes the corresponding homomorphisms.  Certified subgroup inclusions make
every ramification quotient honest, and residual coordinates remain quotient
classes attached to the chosen uniformizer.
\end{definition}
